\section{Threats to Validity}
\label{sec:threats}

\subsection{Construct Validity}

The label denotes generic benign/malware status, not ransomware. SHA-256 identity
and source-index linkage improve traceability but cannot create absent family or
behavior semantics. We therefore avoid ransomware, zero-day, and behavioral
claims.

$H$ is entropy of signed attribution mass, and $V$ is variance across one
attribution vector. Neither is calibrated model uncertainty. The term
``explanation-guided'' refers to how the signals are computed, not proof that they
encode human-understandable malware mechanisms.

Deletion measures model behavior under feature replacement. It does not establish
valid PE edits, causality, robustness to every perturbation, or human usefulness.
EC, CF, and additivity examine different constructs and are not combined into an
OFS.

\subsection{Internal Validity}

Outer test partitions are excluded from query and target selection according to
the retained protocol. Policies share initial sets within fold/seed, limiting
confounding by starting labels. The fixed policy weights and min--max
normalization, however, were not evaluated through an independent outer
sensitivity study. The measured comparison may depend on those design choices.

The oracle reveals existing labels without error. Real analysts can abstain,
disagree, or use different evidence, changing both label quality and cost. Batch
selection has no diversity constraint and may contain redundant samples.

Mean/median baselines are calculated from training data, avoiding direct test
leakage. Yet top-SHAP and random masks are not matched by displacement magnitude,
used-feature status, or manifold distance. These confounders can contribute to
the deletion gap.

\subsection{Statistical Conclusion Validity}

Five seed trajectories are averaged within each of five outer folds. The paired
outer fold is the comparison unit; fold--seed combinations are not 25 independent
test datasets. Five folds support effect-size and direction reporting but limited
formal inference. The readiness snapshot does not provide fold-level vectors in
this checkout, so figures assembled here omit uncertainty bars. The submission
artifact must regenerate them directly from retained fold-level files.

Twenty random masks per $k$ characterize mask variability conditional on each
fitted model. The 25 fidelity models share data and folds, so model-level records
are also dependent. Pairwise rank comparisons share models and cannot be counted
as independent observations in a naive significance test.

Restricted-FPR estimates use small outer test folds. With roughly 300 benign
rows, one false positive changes FPR by about 0.33 percentage points. Results at
1\% and 5\% caps do not validate operation at $10^{-3}$ or $10^{-4}$.

\subsection{External Validity}

The 2,000 rows form a short-period subset recovered from EMBER material and do
not reproduce the full official benchmark protocol. Random row folds do not
enforce temporal, family, source, or campaign separation. Results may change
under newer software, another organization, a different benign distribution,
another feature extractor, or adversarial manipulation. External and temporal
validation are required for a broad journal claim.

LightGBM and TreeSHAP are central to the experiment. The acquisition statistics
can be defined for other additive attributions, but their numeric behavior,
cost, and stability may differ for raw-byte networks, graph models, or sequence
models.

\subsection{Reproducibility Validity}

The readiness report names source commit
\texttt{fed4322f58f5bb55ec47833e857370d076592a32}, locked dataset/manifest hashes,
and replay counts. That commit and the confirmatory result directory are not
available in the cloud checkout used for this manuscript revision. We therefore
attribute numerical values to the retained readiness report and do not describe
this writing pass as an independent rerun. Before submission, the exact source
and artifacts must be synchronized, hashes verified, replay executed in a clean
environment, and all tables/figures regenerated from source records.

Computational replay establishes consistency with stored data and code. It does
not prove external label correctness, data-sharing rights, or generalization.
